\scanpage{33}
\mspage{\handwritten{- 44 -}}
% archive sheet number: 32

chacune des fonctions $\handwritten{f}_k$ soit égale à une série entière sommable par rapport aux $x_j-a_{ij}$. Il existe un nombre $\handwritten{\rho>}0$ tel que, pour tout $\handwritten{x\in\varphi}(I\handwritten{\times}J)$, le polycylindre ouvert de centre $\handwritten{x}$ et de rayons $\handwritten{\rho}$ soit contenu dans un $P_i$ au moins ; car si $\handwritten{\delta}(\handwritten{x})$ désigne la plus grande distance de $\handwritten{x}$ aux frontières de ceux des $P_i$ qui contiennent $\handwritten{x}$, on vérifie aussitôt que $\handwritten{\delta}(\handwritten{x})$ est continue dans $\handwritten{\varphi}(I\handwritten{\times}J)$ et est $\handwritten{>}0$ en tout point, donc atteint son minimum $\handwritten{\rho}$ en un point de $\handwritten{\varphi}(I\handwritten{\times}J)$, ce qui implique $\handwritten{\delta}(\handwritten{x})\handwritten{\geq\rho>}0$ pour tout $\handwritten{x\in\varphi}(I\handwritten{\times}J)$. Par suite, pour tout $\handwritten{x\in\varphi}(I\handwritten{\times}J)$, chacune des fonctions $\handwritten{f}_k$ est égale à une série entière sommable (par rapport aux $z_j-\handwritten{x}_j$) dans le polycylindre ouvert de centre $\handwritten{x}$ et de rayon $\handwritten{\rho}$ (chap.I, \handwritten{\S}1, prop.\handwritten{\unsure{4|14}}). En vertu de la continuité uniforme de $\handwritten{\varphi}$, il existe $\handwritten{\varepsilon>}0$ tel que les relations $\handwritten{|}t-t'\handwritten{|}\handwritten{\leq\varepsilon}$, $\handwritten{|}s-s'\handwritten{|}\handwritten{\leq\varepsilon}$ entraînent $\handwritten{\|\varphi}(t,s)-\handwritten{\varphi}(t',s')\handwritten{\|\leq\frac{\rho}{4}}$. Soit $(t_{\handwritten{\alpha}})$ $(0\handwritten{\leq\alpha\leq}p)$ une suite croissante de points de I tels que $t_0=a$, $t_p=b$ et $t_{\handwritten{\alpha}+1}-t_{\handwritten{\alpha}}\handwritten{\leq\varepsilon}$ pour $0\handwritten{\leq\alpha\leq}p-1$ ; soit de même $(s_{\handwritten{\beta}})$ $(0\handwritten{\leq\beta\leq}q)$ une suite croissante de points de J tels que $s_0=c$, $s_q=d$ et $s_{\handwritten{\beta}+1}-s_{\handwritten{\beta}}\handwritten{\leq\varepsilon}$ pour $0\handwritten{\leq\beta\leq}q-1$. Désignons par $\handwritten{\gamma_\beta}$ le circuit défini par
\[
\begin{aligned}
\handwritten{\gamma_\beta}(t)&=\handwritten{\varphi}(t_{\handwritten{\alpha}},s_{\handwritten{\beta}})+\frac{t-t_{\handwritten{\alpha}}}{t_{\handwritten{\alpha}+1}-t}\bigl(\handwritten{\varphi}(t_{\handwritten{\alpha}+1},s_{\handwritten{\beta}})-\handwritten{\varphi}(t_{\handwritten{\alpha}},s_{\handwritten{\beta}})\bigr)\quad\text{pour }t_{\handwritten{\alpha}}\handwritten{\leq}t\handwritten{\leq}t_{\handwritten{\alpha}+1}\\
&0\handwritten{\leq\alpha\leq}p-1,\quad1\handwritten{\leq\beta\leq}q-1
\end{aligned}
\]
; on posera en outre $\handwritten{\gamma}_0=\handwritten{\gamma}$, $\handwritten{\gamma}_q=\handwritten{\gamma}'$. Tout revient à prouver que l'on a, pour $0\handwritten{\leq\beta\leq}q-1$, $\handwritten{\int_{\gamma_\beta}}(\handwritten{\sum_k f}_k\,dx_k)=$ $=\handwritten{\int_{\gamma_{\beta+1}}}(\handwritten{\sum_k f}_k\,dx_k)$. Or, pour chaque indice $\handwritten{\alpha}$ tel que $0\handwritten{\leq\alpha\leq}p-1$, les points $\handwritten{\gamma_\beta}(t)$ et $\handwritten{\gamma_{\beta+1}}(t)$ pour $t_{\handwritten{\alpha}}\handwritten{\leq}t\handwritten{\leq}t_{\handwritten{\alpha}+1}$ appartiennent tous à un même polycylindre ouvert $Q_{\handwritten{\alpha},\handwritten{\beta}}$ de centre un point de $\handwritten{\varphi}(I\handwritten{\times}J)$ et de rayons $\handwritten{\rho}$ (fig.1) ; on sait d'autre part qu'il existe dans $Q_{\handwritten{\alpha\beta}}$
